\section{Plans, forecasts and revision across decisions}

Economic strategy also requires maintaining or revising an objective as opportunities arrive. The fixed 240-decision sample contains 59 records with forward-looking own-action passages, 36 with prospective opponent capacity or outcome passages and 25 with explicit correction or feedback references. Categories overlap. Fifty-three private fields are absent, all for Opus, and none of these sampled actions is an engine default. These counts describe text in a filtered sample, not population rates of planning or correction.

Eight sampled records state landing policies to buy named properties if encountered. Searching subsequent actual purchase menus yields six without a later named-target opportunity. Grok's G2 d161 policy is executed when it later buys Mediterranean at d256. Its G16 d103 policy to buy Virginia on sight meets a real opportunity at d159 with \$371 against the \$160 list price; the note still favors purchase, but the received action starts an auction. Gemini subsequently acquires Virginia for \$380. The six unexercised policies cannot be scored as failures, and the latter case cannot be treated as a deliberate abandonment when note and action conflict.

The record also contains sustained achievement. Sol at G9 d64 targets three houses on each dark blue as funds arrive, adds houses at d67 and d102 and reaches three each at d163, across global turns 29--60. The G14 eight-house objective is attained over several same-turn decisions. These are realized multi-stage goals, without a measured causal contribution of remembering the initial note. A shorter financing sequence shows why endpoint choice matters. Grok at G20 d99 says it will mortgage Vermont next, first attempts a sale at d100, and mortgages after rejection at d102. It violates a literal next-action prediction while reaching the same-turn cash endpoint recorded in the note.

Opponent assessments usually concern public economic capacity. Grok's G2 d161 forecast that withdrawal leaves New York to another player is realized after two cash-constrained rivals also drop. Such deductions can be useful without requiring inference about private beliefs. Other claims lack a realized test. G19 d220 predicts that rivals will not pay for its survival, but no offer is presented before elimination. A G17 auction forecast correctly anticipates being outbid, while the rival's subsequent jail and liquidation path supplies no building menu to test an accompanying development forecast. The study therefore records particular capacity deductions and action forecasts rather than a calibrated theory-of-mind score.

Revisions are equally concrete. In a G20 auction, Grok's stated \$235 ceiling becomes \$230, followed by a \$232 bid that explicitly acknowledges exceeding the newer limit; it later withdraws without exceeding the original bound. In G17, ceilings rise as the note supplies new rival-resale arguments. These are revisions of stated rules, with no reference valuation proving that the revision improves expected outcomes. They show that agents update and rationalize commitments during interaction, while leaving the value of those changes open.

The full construction-submission ledger gives a separate recurrence endpoint. Among 114 decisions with an uneven-intermediate-operation error, 109 have a later construction submission by that player in the same game, and 63 repeat the error at that next decision's first submission. In 54 error decisions, a later corrective request carries the exact error. Fifty have a subsequent submission, with nine repetitions. Restricting to the first such request per player-game yields 31 episodes, 29 subsequent submissions and two repetitions. Restricting to exposed decisions that end in an applied build yields 44 subsequent submissions and four repetitions. These overlapping populations distinguish exposure to a request, repair of the current response and later recurrence.

Sixty Opus errors follow a refusal that already exhausts its attempt budget, so the configuration receives no later construction-error request there. Most recurrent errors in the all-episode population belong to that condition. Lower recurrence among request-exposed cases cannot be interpreted as a feedback effect because states, targets, response conditions and selection into later builds differ. The positive result is operationally narrower: there are successful repairs and achieved targets, alongside observable recurrence and revision. The evidence does not establish durable learning, retention after delay or adaptation to unfamiliar opponents.

\section{Related work and current context}

NegotiationArena studies multi-turn allocation and price bargaining and experimentally changes behavioral tactics \citep{bianchi2024negotiationarena}. Its controlled outcome comparisons establish a different evidential basis for tactic effectiveness than our observed settlements. The 2026 MERIT study introduces AgoraBench and utility-grounded bargaining measures with feedback and learning interventions \citep{merit2026}. Our cohort instead follows exchanges whose consequences persist in a shared game. It does not import another task's utility or fairness score to value Monopoly portfolios.

AgentMonopoly, inspected as arXiv:2610.29330v2, directly studies model, roster and persona differences in Monopoly \citep{wu2026agentmonopoly}. It measures actions as well as language and permits bluffing while prohibiting public messages contradicting chosen actions. Monopoly and language/action inconsistency are therefore established study topics. Our behavioral account emphasizes repeated counteroffers, management after completion, debt liquidation and creditor paths. FinBoardBench, inspected as arXiv:2605.27896v1, evaluates Cashflow, Acquire and Monopoly and explicitly records forced interventions after response failures \citep{hu2026finboardbench}. We credit that prior distinction between agent and fallback actions. Its described creditor rules sell buildings to the bank, whereas intact buildings in this cohort connect inheritance to later productive portfolios.

The Vending-Bench Arena communication study, inspected as arXiv:2608.14825v3, analyzes 2,583 emails in twenty year-long simulations \citep{li2026arena}. It combines model classification, simulator checks and reasoning-trace review, with full-pipeline replication using two additional judge families. This is substantial prior work on economically grounded communication analysis, not merely unvalidated text labeling. Its external speech-act codebook and conduct setting differ from our explicit permission for public deception. Our case interpretations center financial alternatives, delivered context and explicit correction without estimating a cross-model misalignment rate.

DECE, inspected as arXiv:2610.35773v1, evaluates complete configurations across ten Eurogames and includes matched harness comparisons \citep{zhang2026dece}. It supports the importance of execution conditions but does not supply a causal estimate for the response differences in our cohort. Agent Evaluation Reliability, arXiv:2610.00651v1, formally distinguishes reliable rankings of complete systems from rankings of underlying models \citep{hardy2026reliability}. That distinction is especially relevant when refusal recovery jointly changes the tool schema, retry access and memory.

Process-sensitive evaluation is also established by AgentBoard's progress measurements \citep{ma2024agentboard}. Measurement-First Auditing, arXiv:2610.05830v1, combines provenance, matched controls and scorer validation \citep{yang2026measurement}. We consequently avoid claiming that making a behavioral judgment inspectable validates a general auditing method. Empty Commitments, arXiv:2610.01045v2, separates promises impossible under runtime affordances from failures to enable feasible action \citep{empty_commitments_v2}. Our landing-policy and withdrawal chronologies concern observed opportunities and changed positions rather than its validated ontology. Work on unfaithful explanations further cautions against treating private text as the computation that caused an action \citep{turpin2023unfaithful}. None of these sources determines the upstream cause of a specific response conflict here.

\section{Discussion}

\subsection{The economic unit is a continuing competitive relationship}

The connected finding is that development opportunities are negotiated in relation to other players' ability to finance and inherit them. The offer census, expressed rival-sensitive reasons, first-management outcomes and creditor transfers describe successive parts of that relationship. The social cases give those dependencies explicit content: withholding development cash, preserving an income-producing mortgage option, waiting for a distressed seller, and preventing another player's estate gain. These are concrete forms of strategic conduct observed in a game, even though the cohort does not identify their expected-value optimality.

Three economic states must remain distinct. Ownership completion creates permission to develop. Accessible liquidity and stock determine which plans can execute. Realized rent and estates determine subsequent resources, conditional on movement and survival. Combining these into a single completion or investment-success score would erase the causal order that agents confront. It would also count terminal inherited groups as unused opportunities and system defaults as expressed reserve preferences. The present evidence supplies a descriptive trajectory through those states, with successes, constraints and execution failures all retained.

This is useful beyond the observation that some attractive trades are rejected. The meaningful result is a structured account of contested development: rejection is common where completion favors the current author, bargaining continues where the recipient buys completion, some text expressly weighs rival financing, and the resulting groups have heterogeneous paths into production. The cash differences and changing game states are part of that account. They prevent an isolated causal explanation while making the actual bargaining object more specific than a count of generic agreements.

\subsection{Coverage of the eight research questions}

The user's motivating questions span more than this cohort measures. For strategic reasoning, Q1, we establish selective responses to different exchanges and documented consideration of rival consequences, together with particular legal financial successes and mistakes. We do not measure maximized winning probability. For opponent modeling, Q2, there are checkable public-capacity forecasts and attempts to anticipate responses, without calibrated hidden-belief prediction. For cooperation, negotiation and deception, Q3, the evidence is strongest on complete response behavior and exact forms of bluffing, omission, withdrawal and correction. The conduct is authorized under the public-message instruction; causal persuasion and a fair-surplus optimum remain unmeasured.

For long-horizon planning and memory, Q4, the sustained dark-blue goal and opportunity-conditioned landing policies show both follow-through and execution conflict. The bounded histories and action-only memory condition prevent a general memory-benefit estimate. For adaptation and learning, Q5, ceilings change, retries repair actions and some encoding errors recur. Improvement against a controlled feedback contrast or unfamiliar opponent is not established. For economics and risk, Q6, the study documents actual expenditures, accessible finance, liquidation and receipt exposure. Those measurements do not answer whether \$400 is an optimal price for a property under uncertain continuation.

For emergent and collective behavior, Q7, typed transfers quantify asset concentration and the G20 incident documents attempted coordination over cash. Ordinary dense trading or creditor concentration is not a coalition measure. For evaluation validity and generalization, Q8, the reconstruction checks support specific historical facts, while response and observation defects limit model attribution. The narrow roster and failed readiness comparators leave transfer to stronger or unfamiliar opponents untested. Appendix~\ref{app:coverage} maps all eight questions to observable endpoints and the remaining evidence gaps without declaring the broad agenda complete.

\subsection{The three broader questions}

The distinction between sounding strategic and executing strategy has a mixed, concrete answer here. Some expressed plans become multi-stage development or staged debt payment. Other feasible allocations fail serialization, a planned purchase becomes an auction, and a sale-oriented response becomes bankruptcy. Strategic prose therefore supplies targets and reasons to inspect, not a substitute for successful action. These outcomes belong to the evaluated system, including its orchestration policy, and cannot all be assigned to a model's economic preferences.

The relationship between general and multi-agent capability is unidentified. Four heterogeneous configurations, no independently measured capability axis and unequal effective response conditions cannot estimate that association. Their win totals do not demonstrate either transfer of general intelligence or its failure. The literature's larger cross-game and matched studies address parts of this broader question under their own configurations; they do not fill the missing design in this cohort.

Finally, the records reveal some routes through which mistakes, movement and inheritance enter outcomes, but do not decompose wins into skill, luck and exploitation. A debt-payment witness changes a local endpoint; it does not preserve the historical future after the intervention. An eventual winner's creditor receipts are selected by realized play. Shared seeds likewise do not maintain matched subsequent landings. The observed error and transfer mechanisms are explanatory evidence about specific trajectories, not an allocation of causal credit for victory.

\section{Limitations}

The study's target is one finite historical cohort. Five seed blocks, four cyclic orders and one fixed opponent roster constrain generalization. Decisions and repeated offers are dependent, and the source/refusal boundary is confounded with the first block. We report finite-corpus counts rather than attach decision-level sampling intervals that imply an independent population of plays. Exposure conditioning also selects states already produced by prior actions; it cannot repair the design for a model comparison.

Implementation choices and defects are substantive experimental conditions. Intact creditor inheritance, even-building rules, access to declaration, the history formatter, unequal retry paths and received-response conflicts change what is played and observed. The redemption rounding defect remains preserved rather than repaired or experimentally varied here. Historical replay and legal witnesses share the implemented rules and cannot independently validate them. Finding a defect does not identify the outcome of fixing it.

Semantic categories and incidents were selected retrospectively with LLM research assistance, including family overlap and partial label exposure. The exhaustive mechanical offer population does not make its reason codes objectively correct. The temporal sample is filtered, provider and private channels are unevenly available, and an exact quote can still be a false explanation. We supply neither a human gold standard nor a population deception, intent or planning rate. Recipient delivery establishes opportunity to receive a message, while explicit correction or endorsement supplies narrower textual evidence than latent belief.

Local alternatives use stated endpoints and own-action restrictions. The negative checker cannot certify general bundled-sale impossibility. Current payment feasibility does not establish future survival or optimal continuation. Realized receipts omit unpaid terminal demands and do not estimate incremental returns. Existing reference policies and scenario fixtures do not supply validated economic optima; their readiness defects are summarized separately in the appendix. No new controlled evaluation, game, model call, reference-policy run or synthetic verifier test was performed for this rewrite.

\section{Conclusion}

Across twenty complete games, these configured agents bargain over opportunities whose value depends on rival financing, executable development and creditor succession. The complete offer record shows differentiated responses to development exchanges; contextual transcripts reveal the corresponding pressure, bluffing and corrections. Development outcomes and liquidation evidence then explain why possession, attempted investment and accessible wealth are different economic states. Inheritance concentrates the resulting portfolios in eventual creditors. This connected behavioral account preserves both effective local conduct and consequential execution failures. It supports conclusions about the historical systems and their recorded relationships, while leaving optimal strategy, causal persuasion, durable learning and general capability comparisons for designs that actually measure them.
